%==============================================================================
%  MetaLab Suite --- Problem #2
%  Dynamic Drone Routing Problem (DDRP): theoretical framework
%
%  Compile with:  pdflatex dynamic_drone_routing.tex   (twice for the ToC)
%==============================================================================
\documentclass[11pt,a4paper]{article}

\usepackage[T1]{fontenc}
\usepackage[utf8]{inputenc}
\usepackage{lmodern}
\usepackage[margin=2.5cm]{geometry}
\usepackage{amsmath,amssymb,amsthm,mathtools}
\usepackage{booktabs}
\usepackage{array}
\usepackage{float}
\usepackage{xcolor}
\usepackage{microtype}
\usepackage[hidelinks]{hyperref}

\newtheorem{remark}{Remark}
\DeclareMathOperator*{\argmin}{arg\,min}
\newcommand{\R}{\mathbb{R}}
\newcommand{\E}{\mathbb{E}}

\title{%
  \textbf{The Dynamic Drone Routing Problem}\\[2mm]
  \large Modelling, Objective and Trajectory-Based Metaheuristics}
\author{%
  MetaLab Suite\\[1mm]
  \small Problem \#2 --- Department of Operations Research}
\date{\today}

%==============================================================================
\begin{document}
\maketitle

\begin{abstract}
\noindent
This document develops the theoretical framework for the \emph{Dynamic Drone
Routing Problem} (DDRP), following the drone-delivery vehicle-routing model of
Dorling et al.~(2016). A fleet of unmanned aerial vehicles, launched from a
depot, serves delivery orders that are revealed \emph{in real time}. Each order
carries a payload, and every drone is limited by a strict payload capacity $Q$
and by a battery budget whose consumption scales with both the payload carried
and the distance flown. We give a complete mixed-integer formulation of the
static core, define the objective as the sum of travelled distance and
unserved-order penalties, state the strict constraints, and extend the model to
the online setting through a rolling-horizon decomposition. Designing the
dispatch policy is the task left to the reader.
\end{abstract}

\tableofcontents

%==============================================================================
\section{Introduction}
%==============================================================================

Unmanned aerial vehicles (drones) are increasingly used for last-mile parcel
delivery. Dorling et al.~\cite{dorling2016} introduced and studied the
\emph{Vehicle Routing Problem for drone delivery}: a fleet of multi-rotor
drones, launched from a depot, serves a set of customers under hard limits on
payload weight and battery energy. Their analysis shows that the cost of a
delivery is governed \emph{jointly} by distance and by payload-dependent energy
consumption, so routing and loading decisions are inseparable.

In the \emph{Dynamic} Drone Routing Problem studied here, orders are not known
in advance: they arrive over a planning horizon and must be served as they
appear. A dispatcher must therefore repeatedly reassign drones, trading off
exploitation of the current plan against re-optimisation as new information
arrives. The static problem generalises the classical vehicle-routing problem
and is NP-hard \cite{toth2014}; the online nature makes it harder still
\cite{psaraftis1995}.

%==============================================================================
\section{The static core: sets, parameters and variables}
%==============================================================================

\begin{table}[H]
\centering
\caption{DDRP data and decision variables.}
\label{tab:ddrp}
\begin{tabular}{@{}ll@{}}
\toprule
Symbol & Meaning \\
\midrule
$\mathcal{K}=\{1,\dots,K\}$      & set of drones \\
$\mathcal{I}=\{1,\dots,n\}$      & set of delivery orders (customers) \\
$0$                              & the depot (\emph{departure and return node}) \\
$V=\{0\}\cup\mathcal{I}$         & node set \\
$d_{ij}$                         & travel distance between nodes $i,j$ \\
$w_{i}$                          & payload (parcel weight) of order $i$ \\
$Q$                              & payload capacity of a drone \\
$B$                              & usable battery energy of a drone \\
$a$                              & empty-drone energy consumption per unit distance \\
$b$                              & payload energy sensitivity per unit distance \\
$v$                              & cruise speed (distance per unit time) \\
$\tau$                           & service (hover/landing) time per order \\
$\lambda$                        & penalty for an unserved order \\
\addlinespace
$x_{ijk}\in\{0,1\}$              & $1$ if drone $k$ travels directly from $i$ to $j$ \\
$y_{ik}\in\{0,1\}$               & $1$ if drone $k$ serves order $i$ \\
$q_{ijk}\ge 0$                   & payload carried by drone $k$ on arc $(i,j)$ \\
$t_{ik}\ge 0$                    & auxiliary ordering variable of node $i$ in route $k$ \\
$z_{k}\in\{0,1\}$                & $1$ if drone $k$ is used \\
\bottomrule
\end{tabular}
\end{table}

%==============================================================================
\section{Payload-dependent battery model}
%==============================================================================

A multi-rotor's power requirement increases with its total mass, hence with the
payload it carries. Following the linearisation used by Dorling et al.\
\cite{dorling2016}, the \emph{energy consumed per unit distance} by a drone
carrying a payload $q$ is
\begin{equation}
  e(q) = a + b\,q, \qquad a \ge 0,\ b \ge 0 ,
  \label{eq:energy-rate}
\end{equation}
so that the energy spent traversing arc $(i,j)$ with payload $q$ is
$e(q)\,d_{ij}$. The constant $a$ is the empty-drone consumption and $b\,q$ is
the extra consumption caused by the carried mass. This dependency is what
distinguishes DDRP from a pure distance-minimisation problem: a heavy parcel
carried far consumes disproportionately more battery.

%==============================================================================
\section{Objective function}
%==============================================================================

Let $y_{ik}=\sum_{j\in V} x_{ijk}$ indicate whether drone $k$ serves order $i$.
The dispatcher trades off total travelled distance against the number of orders
left unserved. The objective---\emph{minimise distance and unserved
penalties}---reads
\begin{equation}
  \boxed{\;
  \min \ Z(X)
  = \underbrace{\sum_{k\in\mathcal{K}} \sum_{i\in V} \sum_{j\in V}
      d_{ij}\, x_{ijk}}_{\text{total travelled distance}}
  \;+\; \lambda
    \underbrace{\sum_{i\in\mathcal{I}} \Bigl(1-\sum_{k\in\mathcal{K}} y_{ik}\Bigr)}_{\text{unserved orders}}
  \;}
  \label{eq:objective-ddrp}
\end{equation}
where $\lambda \gg \max_{ij} d_{ij}$ is chosen large enough that serving an
order is always preferable to abandoning it, yet finite so that genuinely
unreachable orders (given the battery budget) are penalised rather than making
the instance infeasible. If energy is to be priced explicitly, as in
\cite{dorling2016}, the objective is augmented with
$c_{e}\sum_{k}\sum_{i,j}(a+b\,q_{ijk})\,d_{ij}\,x_{ijk}$, where $c_{e}$ is the
cost per unit of energy; we keep the two-term form \eqref{eq:objective-ddrp} as
the canonical teaching objective.

%==============================================================================
\section{Constraints}
%==============================================================================

The following constraints are \emph{strict}: no plan may violate them.

\paragraph{(C1) Unique assignment and unserved accounting.}
Each order is served by at most one drone,
\begin{equation}
  \sum_{k\in\mathcal{K}} y_{ik} \ \le \ 1,
  \qquad i \in \mathcal{I}.
  \label{eq:ddrp-assign}
\end{equation}

\paragraph{(C2) Flow conservation and route structure.}
For every served node the in-degree equals the out-degree, and every used drone
leaves and re-enters the depot,
\begin{equation}
  \sum_{j\in V} x_{ijk} = \sum_{j\in V} x_{jik} = y_{ik},
  \quad i\in\mathcal{I},\ k\in\mathcal{K};
  \qquad
  \sum_{j\in\mathcal{I}} x_{0jk} = \sum_{i\in\mathcal{I}} x_{i0k} = z_{k},
  \quad k\in\mathcal{K}.
  \label{eq:ddrp-flow}
\end{equation}

\paragraph{(C3) Payload capacity.}
The payload flowing into a customer equals the payload flowing out plus the
parcel delivered there; the depot loads each drone with exactly the parcels it
will deliver:
\begin{align}
  \sum_{j\in V} q_{ijk} - \sum_{j\in V} q_{jik} &= -\,w_{i}\,y_{ik},
    & i &\in \mathcal{I},\ k\in\mathcal{K},
    \label{eq:ddrp-load-balance}\\
  \sum_{j\in V} q_{0jk} &= \sum_{i\in\mathcal{I}} w_{i}\,y_{ik},
    & k &\in\mathcal{K},
    \label{eq:ddrp-depot-load}
\end{align}
and the capacity limit holds on every arc and for the whole load,
\begin{equation}
  0 \ \le \ q_{ijk} \ \le \ Q\,x_{ijk},
  \quad i,j\in V,\ k\in\mathcal{K};
  \qquad
  \sum_{i\in\mathcal{I}} w_{i}\,y_{ik} \ \le \ Q\,z_{k},
  \quad k\in\mathcal{K}.
  \label{eq:ddrp-capacity}
\end{equation}

\paragraph{(C4) Battery budget.}
The energy consumed along the whole route, with consumption scaling with
payload \eqref{eq:energy-rate} and distance, must not exceed the usable battery
energy,
\begin{equation}
  \sum_{i\in V}\sum_{j\in V} \bigl(a + b\,q_{ijk}\bigr)\, d_{ij}\, x_{ijk}
  \ \le \ B\, z_{k},
  \qquad k\in\mathcal{K}.
  \label{eq:ddrp-battery}
\end{equation}
In practice one reserves a safety margin $B_{\text{res}}$ for wind, hovering and
landing, replacing $B$ by $B-B_{\text{res}}$.

\paragraph{(C5) Subtour elimination.}
Auxiliary ordering variables with a big-$M$ bound forbid disconnected subtours,
\begin{equation}
  t_{ik} - t_{jk} + n\,x_{ijk} \ \le \ n-1,
  \qquad i,j\in\mathcal{I},\ i\neq j,\ k\in\mathcal{K}.
  \label{eq:ddrp-mtz}
\end{equation}

\paragraph{(C6) Release times (dynamic information).}
An order $i$ with release time $r_{i}$ may be served only once it is known. With
$t^{a}_{ik}$ the arrival time at order $i$,
\begin{equation}
  y_{ik} = 1 \ \Rightarrow \ t^{a}_{ik} \ \ge \ r_{i},
  \label{eq:ddrp-release}
\end{equation}
linearised as $t^{a}_{ik} \ge r_{i} - M(1-y_{ik})$. Optional time windows
$[\underline{e}_{i},\overline{e}_{i}]$ extend this to
$\underline{e}_{i}\le t^{a}_{ik}\le \overline{e}_{i}$ when $y_{ik}=1$.

\begin{remark}
Constraints (C1)--(C5) define the static DDRP. Its feasible region is discrete
and its objective \eqref{eq:objective-ddrp} is a linear function of the decision
variables, so the static core is a (NP-hard) mixed-integer program.
\end{remark}

%==============================================================================
\section{The dynamic setting}
%==============================================================================

In DDRP the order set $\mathcal{I}$ is not known a priori. Orders arrive
according to a stochastic process---typically a non-homogeneous Poisson process
with rate $\nu(t)$ over a horizon $[0,T]$---and each order becomes
\emph{actionable} at its release time $r_{i}$. The dispatcher observes the state
\begin{equation}
  S_{t} = \bigl(\text{drone positions},\ \text{remaining batteries},\
          \text{served orders},\ \text{pending orders}\bigr)
\end{equation}
and selects an action: dispatch a drone, insert a customer into an existing
route, or idle. A \emph{policy} $\pi$ maps states to actions, and the Dynamic
DDRP seeks
\begin{equation}
  \min_{\pi} \ \E\Bigl[ \ \sum_{t=0}^{T}
      \Bigl( \textstyle\sum_{k}\sum_{i,j} d_{ij}\,x_{ijk}^{t}
             \;+\; \lambda\, U_{t} \Bigr) \ \Bigr],
  \label{eq:ddrp-dynamic}
\end{equation}
where $U_{t}$ is the number of still-unserved, already-released orders at epoch
$t$, and the expectation is over the unknown arrival process.

A tractable and widely used approach is \emph{rolling-horizon
re-optimisation}: at each decision epoch the dispatcher solves a \emph{static}
instance consisting of all pending orders plus new arrivals, using the drones'
current positions, remaining payload and remaining battery as initial
conditions, and commits to the first portion of the plan before re-optimising
at the next epoch. This turns the online problem into a sequence of offline
DDRP instances, each to be solved by your dispatch policy.

%==============================================================================
\section{Complexity}
%==============================================================================

The static DDRP contains the capacitated vehicle-routing problem as a special
case and is therefore NP-hard \cite{toth2014}; the payload-dependent energy model
\eqref{eq:energy-rate} strengthens the coupling between routing and loading. In
the dynamic setting the orders are revealed online, so no planner can know the
future: this is an online stochastic optimisation problem \cite{psaraftis1995}.
Designing the dispatch policy --- subject to the payload, battery and real-time
constraints above --- is the task left to you.

%==============================================================================
\begin{thebibliography}{9}

\bibitem{dorling2016}
K.~Dorling, J.~Heinrichs, G.~G.~Messier, and S.~Magierowski,
``Vehicle routing problems for drone delivery,''
\emph{IEEE Transactions on Systems, Man, and Cybernetics: Systems},
vol.~47, no.~1, pp.~70--85, 2017 (available online 2016).

\bibitem{toth2014}
P.~Toth and D.~Vigo, Eds.,
\emph{Vehicle Routing: Problems, Methods, and Applications}, 2nd ed.
Philadelphia, PA: SIAM, 2014.

\bibitem{psaraftis1995}
H.~N.~Psaraftis,
``Dynamic vehicle routing: Status and prospects,''
\emph{Annals of Operations Research}, vol.~61, no.~1, pp.~143--164, 1995.

\end{thebibliography}

\end{document}
